\documentclass[11pt,a4paper]{article}
\usepackage[margin=25mm,headheight=15pt]{geometry}
\usepackage{fontspec}
\setmainfont{NotoSerif}[Path=fonts/,Extension=.ttf,UprightFont=*-Regular,
 BoldFont=*-Bold,ItalicFont=*-Italic,BoldItalicFont=*-BoldItalic]
\setsansfont{NotoSans}[Path=fonts/,Extension=.ttf,UprightFont=*-Regular,BoldFont=*-Bold]
\setmonofont{NotoSansMono-Regular.ttf}[Path=fonts/,Scale=0.78]
\usepackage{amsmath,booktabs,tabularx}
\usepackage{graphicx,xcolor,tikz}
\usetikzlibrary{arrows.meta,positioning,fit,backgrounds}
\usepackage{microtype}
\usepackage[numbers,sort&compress]{natbib}
\usepackage{xurl,hyperref,fancyhdr,caption}
\definecolor{ink}{HTML}{17334A}
\definecolor{teal}{HTML}{187F79}
\definecolor{blue}{HTML}{396CA3}
\definecolor{muted}{HTML}{546474}
\hypersetup{colorlinks=true,linkcolor=ink,citecolor=teal,urlcolor=blue,
 pdftitle={P0: Compact Kazakh-Russian Translation with Synthetic Supervision},
 pdfauthor={Ruslan Rakhimov}}
\captionsetup{font=small,labelfont=bf}
\setlength{\parskip}{3pt}
\setlength{\parindent}{1em}
\setlength{\emergencystretch}{2em}
\renewcommand{\arraystretch}{1.14}
\pagestyle{fancy}\fancyhf{}
\fancyhead[L]{\small\sffamily\color{muted}P0 \textbullet\ Kazakh--Russian translation}
\fancyhead[R]{\small\sffamily\color{muted}Technical report / working draft}
\fancyfoot[L]{\small\color{muted}31 October 2026}
\fancyfoot[R]{\small\thepage}
\renewcommand{\headrulewidth}{0.3pt}
\newcommand{\PZeroBLEU}{20.14}
\newcommand{\PZeroChrF}{49.55}
\newcommand{\PZeroChrFPP}{46.48}
\newcommand{\PZeroTER}{69.84}
\newcommand{\KazBLEU}{18.89}
\newcommand{\KazChrF}{48.70}
\newcommand{\KazChrFPP}{45.48}
\newcommand{\KazTER}{71.47}
\newcommand{\DeltaBLEU}{1.25}
\newcommand{\DeltaChrFPP}{1.00}
\newcommand{\DeltaTER}{1.62}
\newcommand{\EvalSeconds}{128.0}
\newcommand{\EvalThroughput}{7.91}

\title{\vspace{-1.3cm}\sffamily\bfseries\color{ink}
 P0: Compact Kazakh--Russian Translation\\with NLLB-Generated Supervision}
\author{Ruslan Rakhimov\\\normalsize Kazakhstan-British Technical University (KBTU)}
\date{\small Working draft for supervisor discussion\\31 October 2026}
\begin{document}
\maketitle\thispagestyle{fancy}
\begin{abstract}
P0 is a 202M-parameter autoregressive Transformer for Kazakh-to-Russian
translation. It was trained from scratch on public parallel data, then
continued on a mixture of original pairs and approximately four million
NLLB-generated translations. A fresh FLORES+ devtest evaluation gives
\PZeroBLEU{} BLEU and \PZeroChrFPP{} chrF++, compared with \KazBLEU{} BLEU
and \KazChrFPP{} chrF++ for rescored KazRush outputs. This report describes
the model, training data and evaluation; recovering the original training
logs remains necessary to quantify the contribution of synthetic data.
\end{abstract}

\section{Introduction}
P0 targets local Kazakh-to-Russian translation at roughly 200M parameters.
The main comparator is \texttt{deepvk/kazRush-kk-ru}, a 197M-parameter
T5-based model trained from random initialization on public parallel
corpora~\cite{kazrush,t5}. P0 uses a custom Transformer encoder--decoder
and additional supervision from NLLB-200 3.3B~\cite{transformer,nllb}.
The practical question is whether this compact system can improve on
KazRush while retaining local inference.

\section{Data and training}
\paragraph{Public data.}
The acquisition pipeline supports OPUS, WMT19, KazParC and TIL. The locally
available filtered corpus contains \textbf{6\,851\,060 parallel pairs}.
Filtering normalizes text, removes markup and invalid pairs, limits sentence
length, and deduplicates pairs; language identification and LaBSE filtering
are optional. Historical per-source counts and enabled filtering stages
have not yet been recovered.

\paragraph{Synthetic data.}
Four million Kazakh sentences were selected from
\texttt{kz-transformers/multidomain-kazakh-dataset}, spanning web, news and
books, and translated with \texttt{facebook/nllb-200-3.3B}.
The filtered archive contains \textbf{3\,999\,605 pairs}. This is forward
translation: a teacher supplies Russian targets for Kazakh monolingual
text~\cite{distillation}. Continued training mixed these pairs with the
original corpus; final mixture proportions and deduplicated size await
the server manifest.

A full scan of the available base and synthetic corpora found zero exact
source or source--reference matches with FLORES+ dev/devtest after NFC
normalization and stripping. Near duplicates and teacher benchmark exposure
were not assessed.

\paragraph{Training history and objective.}
The initial public-data stage reportedly plateaued near \textbf{17 BLEU},
after which the synthetic-data continuation was run. The available best
checkpoint records global step 116\,000. The plateau is an author-reported
observation; its evaluation settings and the per-stage learning curves
require the original logs.

Training uses teacher forcing and label-smoothed cross-entropy, ignoring
padding tokens. Current repository defaults are AdamW with
$\beta=(0.9,0.95)$, weight decay 0.01, label smoothing 0.1, gradient clipping
1.0, and a cosine learning rate from $2\times10^{-4}$ to $5\times10^{-5}$
after 2\,000 warmup steps. Micro-batch size is 16 per GPU with four
accumulation steps; BF16 autocast is enabled. These are current defaults,
not verified historical settings for both stages. The saved weights and
fresh evaluation use FP32.

\section{Model}
P0 has \textbf{202\,016\,256 trainable parameters}, 12 encoder layers,
12 decoder layers, hidden size 512 and FFN size 4096. A shared
32\,000-piece SentencePiece unigram vocabulary~\cite{sentencepiece}
supplies source/target embeddings and the tied output projection.
Sequences are limited to 256 subword tokens.

Blocks use pre-RMSNorm~\cite{rmsnorm}, residual connections, SwiGLU
feed-forward layers~\cite{swiglu} and dropout 0.1.
Encoder self-attention is bidirectional with eight heads. Decoder causal
self-attention uses eight query heads and four KV heads
(GQA)~\cite{gqa}; cross-attention uses eight of each.
Self-attention applies per-head Q/K RMSNorm and RoPE~\cite{rope}.
Cross-attention uses neither. The decoder has a final RMSNorm; the encoder
does not. Figure~\ref{fig:architecture} shows the computation.

\begin{figure}[!ht]\centering
\resizebox{0.94\linewidth}{!}{\definecolor{attentionfill}{RGB}{255,226,187}
\definecolor{ffnfill}{RGB}{194,232,247}
\definecolor{embeddingfill}{RGB}{252,224,225}
\definecolor{normfill}{RGB}{242,244,193}
\definecolor{projectionfill}{RGB}{219,223,239}
\definecolor{outputfill}{RGB}{204,231,207}
\definecolor{blockfill}{RGB}{243,243,244}
\begin{tikzpicture}[
 font=\sffamily\small, flow/.style={-{Latex[length=2mm]},thick,draw=black},
 box/.style={draw=black,thick,rounded corners=1.5mm,fill=white,
 minimum width=5cm,minimum height=0.72cm,align=center},
 embedding/.style={box,fill=embeddingfill},
 attention/.style={box,fill=attentionfill},
 ffn/.style={box,fill=ffnfill},
 note/.style={font=\sffamily\scriptsize,text=black,align=center}]
\node[box] (src) at (0,0) {Kazakh source tokens $x$};
\node[box] (tgt) at (7.1,0) {Shifted Russian target tokens $y_{<t}$};
\node[embedding] (se) at (0,1.05) {Token embedding $E$};
\node[embedding] (te) at (7.1,1.05) {Token embedding $E$ (shared)};
\node[attention,minimum height=1.12cm] (ea) at (0,2.65)
 {Pre-RMSNorm $\rightarrow$ self-attention\\
 8 query / 8 KV heads; bidirectional\\Q/K RMSNorm + RoPE; residual add};
\node[ffn,minimum height=0.9cm] (ef) at (0,4.2)
 {Pre-RMSNorm $\rightarrow$ SwiGLU\\$512 \rightarrow 4096 \rightarrow 512$; residual add};
\node[attention,minimum height=1.12cm] (da) at (7.1,2.65)
 {Pre-RMSNorm $\rightarrow$ causal self-attention\\
 8 query / 4 KV heads (GQA)\\Q/K RMSNorm + RoPE; residual add};
\node[attention,minimum height=0.95cm] (dc) at (7.1,4.2)
 {Pre-RMSNorm $\rightarrow$ cross-attention\\8 query / 8 KV heads; residual add};
\node[ffn,minimum height=0.9cm] (df) at (7.1,5.6)
 {Pre-RMSNorm $\rightarrow$ SwiGLU\\$512 \rightarrow 4096 \rightarrow 512$; residual add};
\begin{scope}[on background layer]
 \node[draw=black,thick,rounded corners=3mm,fill=blockfill,
 fit=(ea)(ef),inner xsep=3mm,inner ysep=4mm] (eblock) {};
 \node[draw=black,thick,rounded corners=3mm,fill=blockfill,
 fit=(da)(dc)(df),inner xsep=3mm,inner ysep=4mm] (dblock) {};
\end{scope}
\node[note,anchor=south west] at (eblock.north west) {ENCODER BLOCK $\times\,12$};
\node[note,anchor=south west] at (dblock.north west) {DECODER BLOCK $\times\,12$};
\node[box] (mem) at (0,6.4) {Encoder memory $H_x$};
\node[box,fill=normfill] (norm) at (7.1,7.5) {Final decoder RMSNorm};
\node[box,fill=projectionfill] (logits) at (7.1,8.6) {Tied projection $h_t E^\top$; 32\,000 logits};
\node[box,fill=outputfill] (out) at (7.1,9.7) {Next Russian token; repeat until EOS};
\draw[flow] (src)--(se);\draw[flow] (se)--(ea);\draw[flow] (ea)--(ef);\draw[flow] (ef)--(mem);
\draw[flow] (tgt)--(te);\draw[flow] (te)--(da);\draw[flow] (da)--(dc);
\draw[flow] (dc)--(df);\draw[flow] (df)--(norm);\draw[flow] (norm)--(logits);
\draw[flow] (logits)--(out);
\draw[flow]
 (mem.east)--(3.55,6.4)--(3.55,4.2)--node[above,note]{K, V}(dc.west);
\node[note,text width=4.8cm] at (0,8)
 {$d_{\mathrm{model}}=512$, head size $64$\\Shared source/target vocabulary\\
 202\,016\,256 trainable parameters\\[2mm]Training: teacher forcing\\
 Inference: beam search with self-attention KV-cache};
\end{tikzpicture}
}
\caption{P0 encoder--decoder. Each block is repeated 12 times;
source embeddings, target embeddings and the output projection share $E$.}
\label{fig:architecture}
\end{figure}

\section{Evaluation and results}
\paragraph{Protocol.}
Fresh P0 evaluation uses FLORES+~\cite{floresplus}, beam size 5,
length-normalization exponent 0.6, a 256-token generation limit and batch
size 16. The same checkpoint is evaluated on 997 dev and 1\,012 devtest
pairs; decoding settings were fixed before the full devtest run.
Scoring uses SacreBLEU 2.6.0~\cite{sacrebleu}: case-sensitive BLEU with
\texttt{13a} tokenization, chrF, chrF++ and TER. All scores use a 0--100
scale; lower TER is better.

\begin{table}[!ht]\centering
\caption{FLORES+ devtest, 1\,012 pairs, common scoring.
KazRush parameter count is from its model card.}
\begin{tabular}{@{}lrrrrr@{}}
\toprule System & Params & BLEU $\uparrow$ & chrF $\uparrow$ & chrF++ $\uparrow$ & TER $\downarrow$\\
\midrule
KazRush KK--RU & 197M & \KazBLEU & \KazChrF & \KazChrFPP & \KazTER\\
\textbf{P0} & 202M & \textbf{\PZeroBLEU} & \textbf{\PZeroChrF} &
\textbf{\PZeroChrFPP} & \textbf{\PZeroTER}\\
\bottomrule
\end{tabular}
\label{tab:results}
\end{table}

P0 improves over the KazRush archive by \textbf{\DeltaBLEU{} BLEU} and
\textbf{\DeltaChrFPP{} chrF++} (Table~\ref{tab:results}).
P0 outputs were generated afresh; KazRush outputs were rescored from a
1\,012-line archive with assumed positional alignment and unavailable
original decoding settings. Paired bootstrap testing (1\,000 resamples,
seed 12\,345) gives $p=0.0010$ for BLEU and chrF++, conditional on this
alignment. The systems differ in training data as well as architecture.

The fresh dev run gives 20.31 BLEU, 49.46 chrF, 46.40 chrF++ and 69.72 TER,
exactly reproducing the checkpoint's stored metrics. These dev scores are
separate from the devtest comparison. Translation errors remain: on devtest
row 3, P0 renders the reference's \emph{эти открытия} as \emph{эти ярости},
changing the meaning despite fluent surrounding text.

\paragraph{Inference and software checks.}
Devtest generation takes \EvalSeconds{} seconds on an Apple M4 Pro
(24 GiB, MPS, FP32), or \EvalThroughput{} sentences/s at batch 16,
excluding model loading and scoring. Separate encoder/decoder ONNX graphs
are available; greedy outputs match native PyTorch on three short examples.
The export has no decoder KV-cache, and mobile performance is unmeasured.
All 29 repository tests pass; detailed test results, metric signatures,
hashes and additional archived benchmarks remain in the accompanying
\texttt{evidence/} and \texttt{generated/} files, with reproduction commands
in \texttt{README.md}.

\section{Conclusion and remaining work}
P0 reaches \PZeroBLEU{} BLEU on FLORES+ devtest, exceeding the rescored
KazRush archive at a similar parameter scale. The next step is to recover
training logs and evaluate the pre-NLLB checkpoint with the same protocol,
then compare public-only and mixed-data continuation at matched training
budgets. This is needed to attribute gains to synthetic supervision.
KazRush should also be regenerated with a pinned revision and documented
decoding settings.

\bibliographystyle{plainnat}
\begingroup\small
\bibliography{references}
\endgroup
\end{document}
